\documentclass[a4paper,11pt]{article}

\usepackage[T1]{fontenc}
\usepackage[utf8]{inputenc}
\usepackage[italian]{babel}
\usepackage{amsmath,amssymb}
\usepackage{booktabs}
\usepackage{array}
\usepackage{geometry}
\usepackage[colorlinks=true,linkcolor=blue,urlcolor=blue]{hyperref}
\usepackage{listings}
\usepackage{xcolor}
\usepackage{microtype}
\usepackage{parskip}

\geometry{a4paper, top=2.5cm, bottom=2.5cm, left=3cm, right=3cm}

\lstset{
  basicstyle=\small\ttfamily,
  breaklines=true,
  keepspaces=true,
  frame=single,
  backgroundcolor=\color{gray!8},
  xleftmargin=0.5em,
  xrightmargin=0.5em,
}

\title{\textbf{Report}\\[0.4em]
\large Schedulazione Green-Aware di Funzioni Serverless\\
tramite Selezione di Varianti Pareto-Ottimale con Esplorazione UCB}
\author{}
\date{}

\begin{document}

\maketitle

% ======================================================================
\section{Motivazione}
% ======================================================================

Il paradigma FaaS su infrastrutture edge introduce una tensione fondamentale tra due
obiettivi in conflitto:

\begin{enumerate}
  \item \textbf{Precisione del risultato} --- l'utente si aspetta output corretti o di
    alta qualità.
  \item \textbf{Consumo energetico} --- il nodo edge ha risorse limitate e opera in un
    contesto in cui il costo ambientale dell'energia varia nel tempo a seconda della
    quota di rinnovabili disponibile nella rete elettrica locale.
\end{enumerate}

Per molte funzioni esistono \textit{varianti approssimate} che producono risultati
leggermente meno precisi ma richiedono meno energia computazionale. Un sistema
green-aware dovrebbe scegliere automaticamente la variante più adeguata al contesto
corrente: privilegiare l'accuratezza quando la rete è alimentata da fonti rinnovabili,
e risparmiare energia quando l'intensità di carbonio è elevata.

Il presente lavoro estende Serverledge con tre meccanismi coordinati:

\begin{enumerate}
  \item \textbf{Selezione di varianti Pareto-ottimale} guidata da un parametro
    $\lambda$ derivato dall'intensità di carbonio della rete elettrica locale
    (API ElectricityMaps).
  \item \textbf{Funzione sigmoide parametrizzata per zona geografica}, con parametri
    pre-calcolati offline da dati storici.
  \item \textbf{Esplorazione UCB (Upper Confidence Bound)} che promuove la valutazione
    di varianti poco osservate, raffinando le stime energetiche con dati reali raccolti
    da InfluxDB.
\end{enumerate}

% ======================================================================
\section{Funzioni Utilizzate}
% ======================================================================

Il sistema è stato valutato su sei funzioni raggruppate in due famiglie:

\begin{itemize}
  \item \textbf{Funzioni numeriche} (PiLeibniz, LnHarmonic, MonteCarlo): output
    quantificabile, errore misurato come scarto relativo dal valore esatto.
  \item \textbf{Funzioni ML} (SentimentAnalysis, ImageClassification, SpamDetection):
    output di classificazione, qualità misurata come accuracy su benchmark di
    riferimento.
\end{itemize}

\noindent\textit{$\star$ = variante approssimata sul fronte di Pareto.}

% -----------------------------
\subsection{PiLeibniz}
% -----------------------------

Stima $\pi$ tramite la serie di Leibniz con $N$ termini:
\[
  \frac{\pi}{4} = \sum_{k=0}^{N-1} \frac{(-1)^k}{2k+1}
\]
Maggiore è $N$, più accurata è la stima ma maggiore è il costo computazionale.
È la funzione con la frontiera di Pareto più ricca (nove varianti che coprono tre ordini
di grandezza in energia ed errore).

\begin{table}[h!]
\centering
\small
\begin{tabular}{lcccc}
\toprule
\textbf{Variante} & \textbf{N termini} & \textbf{Energia (J)} &
\textbf{Errore rel.} & \textbf{Appross.} \\
\midrule
\texttt{base}          & adattivo & 0.007823 & $0.0$              & No \\
\texttt{optimization-py} & adattivo & 0.007318 & $0.0$            & No \\
\texttt{c}             & adattivo & 0.001052 & $0.0$              & No \\
\texttt{n200000} $\star$ & 200\,000 & 0.001534 & $1.0\times10^{-5}$ & Sì \\
\texttt{n100000} $\star$ & 100\,000 & 0.001163 & $2.0\times10^{-5}$ & Sì \\
\texttt{n50000}  $\star$ &  50\,000 & 0.000817 & $4.0\times10^{-5}$ & Sì \\
\texttt{n10000}  $\star$ &  10\,000 & 0.000421 & $2.0\times10^{-4}$ & Sì \\
\texttt{n5000}   $\star$ &   5\,000 & 0.000308 & $4.0\times10^{-4}$ & Sì \\
\texttt{n1000}   $\star$ &   1\,000 & 0.000196 & $2.0\times10^{-3}$ & Sì \\
\bottomrule
\end{tabular}
\end{table}

% -----------------------------
\subsection{LnHarmonic}
% -----------------------------

Stima $\ln(2)$ tramite la serie armonica alternante con $N$ termini:
\[
  \ln(2) = \sum_{k=1}^{N} \frac{(-1)^{k+1}}{k}
         = 1 - \frac{1}{2} + \frac{1}{3} - \frac{1}{4} + \ldots
\]
Struttura analoga a PiLeibniz: il parametro $N$ bilancia accuratezza e consumo
energetico.

\begin{table}[h!]
\centering
\small
\begin{tabular}{lcccc}
\toprule
\textbf{Variante} & \textbf{N termini} & \textbf{Energia (J)} &
\textbf{Errore rel.} & \textbf{Appross.} \\
\midrule
\texttt{base}          & adattivo & 0.007700 & $0.0$              & No \\
\texttt{n200000} $\star$ & 200\,000 & 0.001534 & $3.6\times10^{-6}$ & Sì \\
\texttt{n100000} $\star$ & 100\,000 & 0.001163 & $7.2\times10^{-6}$ & Sì \\
\texttt{n50000}  $\star$ &  50\,000 & 0.000817 & $2.9\times10^{-5}$ & Sì \\
\texttt{n10000}  $\star$ &  10\,000 & 0.000421 & $1.4\times10^{-4}$ & Sì \\
\texttt{n5000}   $\star$ &   5\,000 & 0.000308 & $2.9\times10^{-4}$ & Sì \\
\texttt{n1000}   $\star$ &   1\,000 & 0.000196 & $1.4\times10^{-3}$ & Sì \\
\bottomrule
\end{tabular}
\end{table}

% -----------------------------
\subsection{MonteCarlo}
% -----------------------------

Stima $\pi$ col metodo Monte Carlo: genera $N$ punti casuali in $[0,1]^2$ e conta
quanti cadono nel cerchio unitario. L'errore statistico cala come $O(1/\sqrt{N})$.
Tutte le varianti sono approssimate (l'algoritmo è intrinsecamente stocastico); il
fronte di Pareto copre l'intero set.

\begin{table}[h!]
\centering
\small
\begin{tabular}{lcccc}
\toprule
\textbf{Variante} & \textbf{N campioni} & \textbf{Energia (J)} &
\textbf{Errore rel.} & \textbf{Appross.} \\
\midrule
\texttt{n1000}   $\star$ &   1\,000 & 0.000147 & 0.0520 & Sì \\
\texttt{n5000}   $\star$ &   5\,000 & 0.000912 & 0.0232 & Sì \\
\texttt{n10000}  $\star$ &  10\,000 & 0.002023 & 0.0164 & Sì \\
\texttt{n50000}  $\star$ &  50\,000 & 0.014308 & 0.0074 & Sì \\
\texttt{n100000} $\star$ & 100\,000 & 0.033710 & 0.0052 & Sì \\
\texttt{n500000} $\star$ & 500\,000 & 0.231480 & 0.0023 & Sì \\
\bottomrule
\end{tabular}
\end{table}

% -----------------------------
\subsection{SentimentAnalysis}
% -----------------------------

Classifica un testo in input come \textbf{POSITIVE} o \textbf{NEGATIVE}
(classificazione binaria). Le varianti sono modelli Transformer di dimensioni
decrescenti più un classificatore lessicale (VADER/AFINN) come variante ultra-light.
Qualità stimata su benchmark SST-2.

\begin{table}[h!]
\centering
\footnotesize
\setlength{\tabcolsep}{4pt}
\begin{tabular}{p{2.2cm}p{5.8cm}cccc}
\toprule
\textbf{Variante} & \textbf{Modello} & \textbf{Params} &
\textbf{Energia (J)} & \textbf{Qualità} & \textbf{Appross.} \\
\midrule
\texttt{roberta-large} & siebert/sentiment-roberta-large    & 355 M & 0.0916 & Alta        & No \\
\texttt{distilbert} $\star$ & distilbert-base-uncased-sst-2 &  67 M & 0.0387 & Media       & Sì \\
\texttt{bert-tiny}  $\star$ & mrm8488/bert-tiny-sst2        &   4 M & 0.0113 & Media-bassa & Sì \\
\texttt{vader}      $\star$ & AFINN lexicon (Python puro)   &  ---  & 0.0024 & Bassa       & Sì \\
\bottomrule
\end{tabular}
\end{table}

% -----------------------------
\subsection{ImageClassification}
% -----------------------------

Classifica un'immagine in input (Base64) tra 1000 categorie ImageNet, restituendo la
categoria predetta e il confidence score. Le varianti sono architetture Transformer/CNN
di dimensioni decrescenti. Qualità indicata come Top-1 accuracy su ImageNet.

\begin{table}[h!]
\centering
\footnotesize
\setlength{\tabcolsep}{4pt}
\begin{tabular}{p{2.8cm}p{5.5cm}cccc}
\toprule
\textbf{Variante} & \textbf{Modello} & \textbf{Params} &
\textbf{Energia (J)} & \textbf{Qualità} & \textbf{Appross.} \\
\midrule
\texttt{vit-base}          & google/vit-base-patch16-224       & 86 M  & 1.4553 & Alta        & No \\
\texttt{efficientnet-b0} $\star$ & google/efficientnet-b0      & 5.3 M & 0.4180 & Media       & Sì \\
\texttt{mobilenet-v2}    $\star$ & google/mobilenet\_v2\_1.0\_224 & 3.4 M & 0.2289 & Media-bassa & Sì \\
\texttt{mobilenet-v2-tiny} $\star$ & google/mobilenet\_v2\_0.35\_224 & 1.7 M & 0.0770 & Bassa  & Sì \\
\bottomrule
\end{tabular}
\end{table}

% -----------------------------
\subsection{SpamDetection}
% -----------------------------

Classifica un messaggio (email, SMS) come \textbf{SPAM} o \textbf{HAM}
(classificazione binaria). Le varianti spaziano da modelli Transformer fine-tuned per
spam detection fino a un classificatore euristico basato su keyword. Qualità stimata
su corpus misti (Enron + SMS Spam Collection + phishing samples).

\begin{table}[h!]
\centering
\footnotesize
\setlength{\tabcolsep}{4pt}
\begin{tabular}{p{2.2cm}p{6.0cm}cccc}
\toprule
\textbf{Variante} & \textbf{Modello} & \textbf{Params} &
\textbf{Energia (J)} & \textbf{Qualità} & \textbf{Appross.} \\
\midrule
\texttt{roberta-spam}  & mshenoda/roberta-spam                         & 125 M & 0.0850 & Alta        & No \\
\texttt{distilbert} $\star$ & Falconsai/spam\_classification            &  67 M & 0.0350 & Media       & Sì \\
\texttt{bert-tiny}  $\star$ & mrm8488/bert-tiny-finetuned-sms-spam     & 4.4 M & 0.0100 & Media-bassa & Sì \\
\texttt{keyword}    $\star$ & Keyword heuristic (Python puro)           &  ---  & 0.0020 & Bassa       & Sì \\
\bottomrule
\end{tabular}
\end{table}

% ======================================================================
\section{Esplorazione UCB}
% ======================================================================

\subsection{Motivazione}

Il profilo energetico offline di ogni variante (\texttt{invocation\_joule} nel JSON di
configurazione) è ottenuto eseguendo ripetutamente le varianti su un hardware di
riferimento e registrando l'energia media consumata.
Questo approccio ha due limitazioni concrete che hanno reso necessario il meccanismo
di esplorazione:

\begin{enumerate}
  \item \textbf{Dipendenza dall'hardware}: i profili energetici misurati sull'hardware
    di sviluppo non corrispondono necessariamente ai consumi reali su un nodo edge con
    CPU o workload diversi. I valori offline potrebbero sovra- o sotto-stimare
    sistematicamente il consumo reale, portando lo scheduler a preferire stabilmente
    una variante nominalmente ``economica'' che in produzione è invece costosa --- e
    viceversa.

  \item \textbf{Scalabilità della profilazione}: per ogni nuova variante sarebbe
    necessario ripetere una profilazione offline esaustiva prima di poterla deployare.
    Questo crea un collo di bottiglia nello sviluppo e impedisce l'aggiunta rapida di
    nuove varianti al sistema.
\end{enumerate}

Indipendentemente da queste due limitazioni principali, vale la pena notare che il
consumo energetico non è stazionario --- dipende dal carico del sistema, dalla
temperatura e dallo stato della cache del container. Una stima offline non cattura
questa variabilità.

Il meccanismo UCB (Upper Confidence Bound) affronta queste problematiche con il
principio di \textbf{ottimismo sotto incertezza} (\textit{Optimism in the Face of
Uncertainty}): per varianti poco osservate, il sistema assume \textit{ottimisticamente}
che il loro consumo energetico reale potrebbe essere inferiore alla stima corrente.
Questo incentiva l'esplorazione di varianti sotto-campionate, raccogliendo dati reali
da InfluxDB e convergendo progressivamente a stime accurate senza richiedere alcuna
profilazione offline preventiva.

\subsection{Come vengono acquisiti i dati}

Il flusso di acquisizione è il seguente:

\begin{lstlisting}
Invocazione
    |
    v
Esecuzione container
    |
    v
Lettura RAPL / perf-stat  (energia CPU durante l'esecuzione)
    |
    v
Write su InfluxDB
    measurement : energy_sample
    tag         : variant_id, function_name
    fields      : invocation_joule, cold_start_joule, timestamp
    |
    v
EMA update su etcd  (alpha = 0.2)
    E_new = 0.2 * E_measured + 0.8 * E_old
\end{lstlisting}

Ogni invocazione produce quindi due effetti persistenti:

\begin{itemize}
  \item \textbf{InfluxDB} accumula la serie storica delle misurazioni per variante; da
    essa si estraggono $\sigma$ e $n$ utilizzati nel termine di esplorazione UCB.
  \item \textbf{etcd} mantiene la stima EMA corrente: è il valore $\mu$ primario nella
    formula LCB, nonché il fallback unico quando InfluxDB non è disponibile o ha
    campioni insufficienti.
\end{itemize}

\subsubsection*{Formula LCB (Lower Confidence Bound)}

Per la \textit{minimizzazione} energetica si usa il limite inferiore di confidenza
(equivalente all'UCB classico applicato alla quantità negata $-E$):

\[
  \text{LCB}_i = \mu^{\text{EMA}}_i
               - \beta \cdot \frac{\sigma^{\text{influx}}_i}{\sqrt{n_i}}
\]

dove:

\begin{table}[h!]
\centering
\small
\begin{tabular}{lp{10cm}}
\toprule
\textbf{Simbolo} & \textbf{Significato} \\
\midrule
$\mu^{\text{EMA}}_i$       & Stima EMA dell'energia della variante $i$
                              (da etcd, $\alpha = 0.2$) \\
$\sigma^{\text{influx}}_i$ & Deviazione standard calcolata sui campioni InfluxDB
                              (ultimi 24 h) \\
$n_i$                      & Numero di campioni InfluxDB disponibili \\
$\beta$                    & Parametro di esplorazione (configurabile, default 1.0) \\
\bottomrule
\end{tabular}
\end{table}

\medskip
\noindent\textbf{Perché EMA come $\mu$?}
L'EMA è un estimatore ricorrente con fattore di smorzamento $\alpha = 0.2$: pondera di
più le misurazioni recenti e dimentica progressivamente quelle vecchie. Questo lo rende
più robusto rispetto a una media aritmetica su 24 h quando il consumo energetico del
nodo varia nel tempo (temperatura, carico CPU, stato cache).

\medskip
\noindent\textbf{Proprietà:}
\begin{itemize}
  \item $n_i$ grande (variante ben esplorata):
    $\sigma_i / \sqrt{n_i} \to 0 \;\Rightarrow\;
    \text{LCB}_i \to \mu^{\text{EMA}}_i$ \emph{(sfruttamento)}
  \item $n_i$ piccolo (variante sotto-esplorata):
    correzione ampia $\Rightarrow \text{LCB}_i \ll \mu^{\text{EMA}}_i$
    \emph{(esplorazione ottimistica)}
  \item $\beta = 0$: UCB disabilitato, si usa solo l'EMA
  \item $\beta \to \infty$: massima esplorazione
\end{itemize}

\medskip
\noindent\textbf{Fallback:} se InfluxDB non è raggiungibile, oppure la variante ha
meno di $\texttt{min\_samples}$ (default: 5) campioni --- soglia sotto la quale
$\sigma$ sarebbe inaffidabile --- lo scheduler usa l'EMA da etcd senza correzione.

% ======================================================================
\section{Scheduling Green-Aware}
% ======================================================================

\subsection{Visione ad alto livello}

Quando arriva una richiesta di invocazione per una funzione logica
(es.\ \texttt{SpamDetection}), lo scheduler esegue i seguenti passi:

\begin{lstlisting}
1. Recupera la lista di varianti fisiche dal JSON (etcd)
2. Per ogni variante: ottieni energia stimata
       --> recupera mu_EMA da etcd
       --> se UCB attivo e n_i >= min_samples: LCB = mu_EMA - beta*sigma_influx/sqrt(n)
       --> altrimenti: usa mu_EMA direttamente (nessuna correzione)
3. Calcola la frontiera di Pareto bi-obiettivo (energia (min), errore (min))
4. Calcola lambda dalla carbon intensity corrente della zona del nodo
5. Scalarizza la frontiera con lambda -> seleziona la variante con score minimo
6. Invoca la variante selezionata; registra l'energia misurata
\end{lstlisting}

L'obiettivo è scegliere la variante che realizza il miglior compromesso tra qualità del
risultato e costo energetico, nel contesto energetico corrente del nodo edge.

\subsection{Il parametro \texorpdfstring{$\lambda$}{lambda}}

\subsubsection*{Frontiera di Pareto}

Dati $n$ varianti, il sistema calcola la frontiera di Pareto bi-obiettivo
(energia, errore), entrambi da minimizzare:

\begin{enumerate}
  \item Ordina le varianti per energia crescente; a parità, per errore crescente.
  \item Scansione singola: una variante entra nella frontiera se e solo se il suo
    errore è strettamente minore del minimo errore dei punti precedenti.
\end{enumerate}

L'algoritmo è $O(n \log n)$ totale e identifica tutti e soli i punti non dominati:
nessuna variante nella frontiera è ``peggiore su entrambi gli assi'' rispetto a
un'altra.

\medskip
\noindent\textbf{Esempio (PiLeibniz, varianti selezionate):}

\begin{table}[h!]
\centering
\small
\begin{tabular}{lccc}
\toprule
\textbf{Variante} & \textbf{Energia (J)} & \textbf{Errore} & \textbf{Pareto?} \\
\midrule
\texttt{n1000}          & 0.000196 & $2.0\times10^{-3}$ & $\checkmark$ \\
\texttt{n5000}          & 0.000308 & $4.0\times10^{-4}$ & $\checkmark$ \\
\texttt{n10000}         & 0.000421 & $2.0\times10^{-4}$ & $\checkmark$ \\
\texttt{n50000}         & 0.000817 & $4.0\times10^{-5}$ & $\checkmark$ \\
\texttt{n100000}        & 0.001163 & $2.0\times10^{-5}$ & $\checkmark$ \\
\texttt{n200000}        & 0.001534 & $1.0\times10^{-5}$ & $\checkmark$ \\
\texttt{c}              & 0.001052 & $0.0$              & $\checkmark$ \\
\texttt{optimization-py} & 0.007318 & $0.0$             & $\times$ dominata da \texttt{c} \\
\texttt{base}           & 0.007823 & $0.0$              & $\times$ dominata da \texttt{c} \\
\bottomrule
\end{tabular}
\end{table}

\subsubsection*{Scalarizzazione con \texorpdfstring{$\lambda$}{lambda}}

Data la frontiera di Pareto, ogni variante $i$ riceve un punteggio:

\[
  \text{score}(i) = (1 - \lambda) \cdot \hat{E}_i + \lambda \cdot \hat{\varepsilon}_i
\]

dove $\hat{E}_i$ e $\hat{\varepsilon}_i$ sono i valori di energia ed errore normalizzati
nell'intervallo $[0,1]$ rispetto al range della frontiera. La variante con il
\textbf{punteggio minimo} viene selezionata.

\medskip
\noindent\textbf{Cosa rappresenta $\lambda$:}

$\lambda \in [0, 1]$ è un \textit{peso di qualità} che esprime quanto privilegiare
l'accuratezza del risultato rispetto all'efficienza energetica nell'istante corrente.

\begin{itemize}
  \item $\lambda \to 0$ (rete sporca, alta CI): il termine dominante è
    $(1-\lambda)\hat{E}_i$ $\to$ il sistema seleziona la variante più efficiente
    energeticamente.
  \item $\lambda \to 1$ (rete pulita, bassa CI): il termine dominante è
    $\lambda\hat{\varepsilon}_i$ $\to$ il sistema seleziona la variante più accurata
    (errore minimo).
  \item $\lambda = 0.5$: selezione al ``ginocchio'' della frontiera, bilanciando i due
    obiettivi.
\end{itemize}

$\lambda$ non è un parametro configurato manualmente: viene calcolato automaticamente
in tempo reale dalla carbon intensity della zona geografica del nodo edge.

\subsection{Calcolo di \texorpdfstring{$\lambda$}{lambda} dalla Carbon Intensity}

La \textbf{carbon intensity} (CI, misurata in gCO$_2$eq/kWh) quantifica quante
emissioni di CO$_2$ equivalente vengono prodotte per ogni kWh di energia consumata
dalla rete elettrica locale. Viene ottenuta in tempo reale dall'API ElectricityMaps
(con cache di 15 minuti).

$\lambda$ è derivato tramite una \textbf{funzione sigmoide logistica} centrata sui
valori storici della zona:

\[
  \lambda = \frac{1}{1 + \exp\!\left(k \cdot \dfrac{CI - \mu_z}{\sigma_z}\right)}
\]

dove:

\begin{table}[h!]
\centering
\small
\begin{tabular}{lp{10cm}}
\toprule
\textbf{Simbolo} & \textbf{Significato} \\
\midrule
$CI$      & Carbon intensity corrente (gCO$_2$eq/kWh) \\
$\mu_z$   & Media storica della CI per la zona $z$ \\
$\sigma_z$ & Parametro di scala:
             $\sigma_z = (P_{95} - P_5) \big/ \!\left(2\ln\tfrac{1-\alpha}{\alpha}\right)$,
             con $\alpha=0.1$ \\
$k$       & Pendenza della sigmoide (default 0.45) \\
\bottomrule
\end{tabular}
\end{table}

La formula per $\sigma_z$ assicura che $\lambda(P_{95}) = \alpha = 0.1$ e
$\lambda(P_5) = 1 - \alpha = 0.9$: ai valori estremi della distribuzione storica,
$\lambda$ raggiunge i valori estremi dell'intervallo, massimizzando la sensibilità del
sistema nelle condizioni reali attese per quella zona.

\subsection{Dipendenza dalla zona geografica}

La sigmoide è \textbf{parametrizzata per zona} perché la distribuzione storica della CI
varia enormemente tra Paesi con mix energetici diversi.

Se si usasse una sigmoide globale con $\mu$ unica, sarebbe completamente insensibile
alle variazioni locali: in Francia (CI media $\approx$ 14 gCO$_2$/kWh), anche i picchi
locali non superano mai la media globale e $\lambda$ sarebbe sempre vicino a 1
indipendentemente dal contesto. In Polonia (CI media $\approx$ 499 gCO$_2$/kWh),
$\lambda$ sarebbe sempre vicino a 0.

Usando $\mu$ e $\sigma$ calibrati sulla distribuzione storica \textit{di quella zona},
la sigmoide risponde alle oscillazioni locali: quando la CI di una zona è sopra la sua
media storica, $\lambda$ scende (si risparmia energia); quando è sotto la media,
$\lambda$ sale (si usa la variante più precisa). Il sistema si adatta all'energia
disponibile nel preciso contesto geografico del nodo, non a una scala globale
arbitraria.

I parametri sono pre-calcolati offline da dati storici scaricati da ElectricityMaps e
salvati in \texttt{zone\_ci\_params.csv}. Le cinque zone scelte coprono l'intero spettro
della carbon intensity europea --- da una rete quasi completamente idroelettrica
(Norvegia) o nucleare (Francia) a una rete a carbone dominante (Polonia) ---
garantendo che il sistema venga validato su regimi di $\lambda$ bassi, medi e alti:

\begin{table}[h!]
\centering
\small
\begin{tabular}{lclp{6.5cm}}
\toprule
\textbf{Zona} & $\boldsymbol{\mu}$ \textbf{(gCO$_2$/kWh)} &
$\boldsymbol{\sigma}$ & \textbf{Caratteristica rete} \\
\midrule
NO & 6   & $\sim$1  & Idroelettrico dominante $\rightarrow$ quasi sempre $\lambda \approx 1$ \\
FR & 14  & $\sim$4  & Nucleare dominante $\rightarrow$ $\lambda$ solitamente alto \\
ES & 95  & $\sim$14 & Mix rinnovabile/gas $\rightarrow$ $\lambda$ variabile \\
DE & 278 & $\sim$41 & Mix carbone/rinnovabili $\rightarrow$ $\lambda$ mediamente basso \\
PL & 499 & $\sim$35 & Carbone dominante $\rightarrow$ quasi sempre $\lambda \approx 0$ \\
\bottomrule
\end{tabular}
\end{table}

% ======================================================================
\section{Bozza degli Esperimenti}
% ======================================================================

Oltre agli esperimenti descritti di seguito, la sperimentazione include alcune analisi
preliminari che ne costituiscono il presupposto: (i) la profilazione offline del consumo
energetico di ogni variante tramite invocazioni ripetute con lettura da InfluxDB;
(ii) la calibrazione della sigmoide per zona geografica a partire dai dati storici di
carbon intensity scaricati da ElectricityMaps.

\subsection{Exp~1 --- Frontiera di Pareto e selezione per
  \texorpdfstring{$\lambda$}{lambda}}

\textbf{Obiettivo:} per ogni famiglia di funzioni, visualizzare la frontiera di Pareto
(energia vs.\ errore/qualità) e mostrare quale variante viene selezionata al variare di
$\lambda$, evidenziando come la scalarizzazione si sposta lungo la frontiera in funzione
del peso attribuito all'efficienza energetica rispetto all'accuratezza.

\medskip
\textbf{Metodologia:}
\begin{itemize}
  \item Per ogni funzione, calcolare la frontiera di Pareto bi-obiettivo
    (energia, errore).
  \item Spazzare $\lambda \in [0, 1]$ a passi di 0.1 e registrare quale variante viene
    selezionata dallo scheduler in assenza di rumore (stime energetiche = valori
    offline).
  \item Visualizzare: scatter plot energia vs.\ errore con fronte evidenziato, e un
    grafico ``variante selezionata vs.\ $\lambda$'' che mostra la transizione lungo il
    fronte.
\end{itemize}

\textbf{Risultato atteso:} la variante selezionata si sposta monotonicamente dalla più
efficiente ($\lambda = 0$) alla più accurata ($\lambda = 1$); le transizioni si
collocano ai valori di $\lambda$ in cui lo score delle varianti si incrociano.

\subsection{Exp~2 --- Comportamento dello scheduling al variare della CI}

\textbf{Obiettivo:} verificare che il sistema selezioni varianti diverse in funzione
della carbon intensity della zona, con una transizione continua dalla variante più
accurata (CI bassa) a quella più efficiente (CI alta) che rispecchi il comportamento
atteso della sigmoide calibrata per zona.

\medskip
\textbf{Metodologia:}
\begin{itemize}
  \item Per ogni zona (NO, FR, ES, DE, PL) e per funzioni rappresentative di entrambe
    le famiglie (es.\ \texttt{PiLeibniz}, \texttt{SpamDetection}):
    \begin{itemize}
      \item Invocare lo scheduler con la CI storica mediana di quella zona.
      \item Registrare la variante selezionata e il valore di $\lambda$.
    \end{itemize}
  \item Mettere in risalto il rapporto fra accuratezza del risultato e consumo
    energetico in funzione dei valori di CI.
\end{itemize}

\textbf{Risultato atteso:}
\begin{itemize}
  \item NO (CI $\approx$ 6) $\;\rightarrow\;$ $\lambda \approx 1$
    $\;\rightarrow\;$ variante più accurata
  \item FR (CI $\approx$ 14) $\;\rightarrow\;$ $\lambda$ alto
    $\;\rightarrow\;$ variante heavy/medium
  \item ES (CI $\approx$ 95) $\;\rightarrow\;$ $\lambda$ medio
    $\;\rightarrow\;$ variante di compromesso
  \item DE (CI $\approx$ 278) $\;\rightarrow\;$ $\lambda$ basso
    $\;\rightarrow\;$ variante leggera
  \item PL (CI $\approx$ 499) $\;\rightarrow\;$ $\lambda \approx 0$
    $\;\rightarrow\;$ variante più efficiente
\end{itemize}

\subsection{Exp~3 --- Risparmio energetico vs.\ baseline always-accurate}

\textbf{Obiettivo:} quantificare il risparmio energetico dello scheduler green-aware
(in termini di CI) rispetto a uno scheduler che seleziona sempre la variante più
accurata, su un workload distribuito su più zone con carbon intensity diverse.

\medskip
\textbf{Metodologia:}
\begin{itemize}
  \item Workload: invocazioni distribuite su tutte e 6 le funzioni, con distribuzione
    uniforme delle zone (NO, FR, ES, DE, PL).
  \item Strategie a confronto:
    \begin{enumerate}
      \item Baseline always-accurate: seleziona sempre la variante con
        \texttt{is\_approximate: false} più precisa, indipendentemente dalla CI.
      \item Scheduler green-aware (questo lavoro).
    \end{enumerate}
  \item Metriche: energia totale (J), risparmio percentuale $\Delta E\%$,
    errore/accuracy medio.
\end{itemize}

\textbf{Risultato atteso:} lo scheduler produce un risparmio energetico significativo
rispetto alla baseline, in particolare nelle zone con CI elevata (DE, PL), mantenendo
qualità alta nelle zone pulite (NO, FR).

\subsection{Exp~4 --- Convergenza EMA con stime iniziali errate}

\textbf{Obiettivo:} verificare che il meccanismo EMA+LCB corregga autonomamente stime
energetiche iniziali errate attraverso le misurazioni reali, mettendo in risalto la
prima fase di esplorazione.

\medskip
\textbf{Metodologia:}
\begin{itemize}
  \item Funzione: \texttt{PiLeibniz}.
  \item Setup: alterare artificialmente \texttt{invocation\_joule} di una variante
    approssimata nel JSON portandolo a un valore significativamente gonfiato rispetto
    al reale.
  \item Forzare un numero iniziale di campioni InfluxDB per attivare il termine LCB,
    quindi eseguire 100 invocazioni di scheduling con $\lambda$ basso (zona DE ---
    preferisce efficienza), registrando ad ogni step: variante selezionata,
    $\mu_{\text{EMA}}$, LCB della variante perturbata.
\end{itemize}

\medskip
\noindent\textbf{Come mostrare la convergenza:}\\
Tracciare su un unico grafico (asse X = invocazioni della variante perturbata):
\begin{itemize}
  \item $\mu_{\text{EMA}}$ che decae esponenzialmente dal valore gonfiato verso il
    valore reale
  \item Soglia corrispondente alla variante concorrente --- quando $\mu_{\text{EMA}}$
    scende sotto questa soglia, la variante perturbata diventa di nuovo competitiva
  \item Punti colorati sulla timeline: quale variante viene selezionata dallo scheduler
\end{itemize}

La convergenza è descrivibile analiticamente: con $\alpha = 0.2$, valore reale $E^*$,
valore iniziale $E_0$:
\[
  \mu^{\text{EMA}}(t) = E^* + (E_0 - E^*) \cdot (1-\alpha)^t
\]

\textbf{Risultato atteso:} il meccanismo LCB rende la variante perturbata attrattiva
in fase esplorativa; man mano che si accumulano misurazioni reali, l'EMA converge al
valore reale e la selezione diventa stabile e corretta.

\subsection{Exp~5 --- Sensibilità al parametro
  \texorpdfstring{$\beta$}{beta}}

\textbf{Obiettivo:} $\beta = 0$ corrisponde al puro sfruttamento (nessuna esplorazione);
valori maggiori aumentano la robustezza a stime iniziali errate a scapito di qualche
invocazione esplorativa in più. Questo esperimento confronta questi regimi.

\medskip
\textbf{Metodologia:}
\begin{itemize}
  \item Funzione: \texttt{PiLeibniz} (fronte di Pareto ricco, 9 varianti).
  \item Setup: stime energetiche iniziali volutamente imprecise (pochi campioni
    offline).
  \item $\beta \in \{0,\; 0.5,\; 1.0,\; 2.0\}$.
  \item Per ciascun $\beta$: 200 invocazioni, registrando variante selezionata ad ogni
    step, numero di invocazioni su varianti sub-ottimali, e invocazione di convergenza.
\end{itemize}

\textbf{Risultato atteso:}
\begin{itemize}
  \item $\beta = 0$: puro sfruttamento --- se la stima iniziale è errata, non converge
    mai.
  \item $\beta = 0.5$--$1.0$: convergenza in 15--30 invocazioni con overhead
    esplorativo limitato.
  \item $\beta = 2.0$: esplora molto di più, convergenza più lenta ma più robusta a
    outlier.
\end{itemize}

\subsection{Metriche di valutazione}

\begin{table}[h!]
\centering
\small
\begin{tabular}{lp{10cm}}
\toprule
\textbf{Metrica} & \textbf{Descrizione} \\
\midrule
$E_{\text{tot}}$          & Energia totale consumata (J) in $N$ invocazioni \\
$\varepsilon_{\text{med}}$ & Errore relativo mediano dell'output (funzioni numeriche) \\
$A_{\text{med}}$           & Accuracy media delle classificazioni (funzioni ML) \\
$\Delta E\%$               & Risparmio energetico vs.\ baseline always-accurate:
                             $(E_\text{base} - E_\text{sched})/E_\text{base}$ \\
$n_{\text{exp}}$           & Numero di invocazioni esplorative
                             (UCB attivo, $n_i < \texttt{min\_samples}$) \\
$t_{\text{conv}}$          & Invocazioni fino alla convergenza UCB
                             (variante ottimale stabile) \\
$\lambda_{\text{med}}$     & Valore mediano di $\lambda$ per zona
                             (verifica calibrazione sigmoide) \\
\bottomrule
\end{tabular}
\end{table}

\end{document}
